The first atomic bomb was dropped on Hiroshima, Japan.
The entire landscape of the city was destroyed. There were short term affects
and long term affects. There were also psychological issues that were increasing
among the survivors. The environmental affects are still being dealt with today,
according to some scientists. High radiation exposure lead to:

\begin{itemize}
  \item Blindness
  \item Hair Loss
  \item Internal and External Bleeding
  \item Vomiting
  \item Etc \ldots
\end{itemize}

Some survivors even experienced bone marrow deterioration, inflammation in the
intestines, cancer, which included lymphoma, breast cancer, and malignant
tumors. Leukemia in children quadrupled. The after effects of the radiation
caused over $1,900$ deaths. People who were near the center began to show signs
of radiation illnesses within two to three days of exposure. But, the people who
were out didn't see the signs until one to four weeks. Women who were pregnant
during the time had an increase chance of birth defects, death, and
miscarriages of their infants. The rate of mental retardation increased
dramatically.

A large number of people died from secondary disease as a result of lowerd
resistance. The inability to treatment and facilities also led to many of the
secondary deaths. It's estimated that if Hiroshima had good health care, then
ten to twenty percent would have survived.

Immediately after the war, the US government and the Japanese government created
the \bf{Atomic Bomb Casualty Commission}. According to the \bf{ABCC}, cancer
started to double with radiation exposure. Over the years it has been documented
that one of the side effects was severe cataracts in those exposed to the
radiation.

The major problem that is still occurring up until today is increase in cancer
and other diseases. Also, the increase of mental instability, birth defects, etc
started to rise ever since the bomb. All of this could have been prevented if
Hiroshima clean and had a good medical system. So, that's my solution to these
ongoing problems. I don't know exactly how to do it, but Japan or some other
country could help Japan build a medical system in Hiroshima. Build multiple
hospitals scattered around Japan, not just Hiroshima. They should also think
about constructing therapy buildings to help those who are mentally unstable.
They should also try to reach out to the $127,000$ survivors to see if they need
any financial, mental, physical help.

\newpage

\begin{center}
  \subsubsection{Citations}
\end{center}

McCamy, Laura. "Here's What Hiroshima Looks like Today - and How the Effects
of the Bombing Still Linger." \it{Business Insider, Business Insider}, 3 Aug.
2018,

Jordan, Bertrand R. "The Hiroshima/Nagasaki Survivor Studies: Discrepancies
between Results and General Perception." \it{Genetics}, Genetics Society of
America, Aug. 2016, https://www.ncbi.nlm.nih.gov/pmc/articles/PMC4981260/. 

Listwa, Dan, and Dan Listwa. “Hiroshima and Nagasaki: The Long Term Health
Effects.” \it{K=1 Project}, 9 Aug. 2012, \\
https://k1project.columbia.edu/news/hiroshima-and-nagasaki. 

